\documentclass[a4paper]{article}

\usepackage[T1]{fontenc}
\usepackage[utf8]{inputenc}
\usepackage[ngerman]{babel}
\usepackage{amsmath,amssymb,amsthm}
\usepackage{tikz}
\usetikzlibrary{patterns}
\usepackage{float}
\usepackage{enumitem}
\usepackage[colorlinks=true,linkcolor=red,citecolor=red,urlcolor=red]{hyperref}

% --- Umgebungen (Stil des Kurzskripts: fettes Label, kursiver Text) ---------
\newtheoremstyle{kurzskript}%
  {\medskipamount}{\medskipamount}% Abstand oben / unten
  {\itshape}{}%                      Textschrift, Einzug
  {\bfseries}{}{ }%                  Labelschrift, Punkt, Abstand
  {\thmname{#1}\thmnumber{ #2}\thmnote{ \bfseries(#3)}}
\theoremstyle{kurzskript}
\newtheorem{satz}{Satz}[section]             % Nummerierung 1.1, 1.2, ...
\newtheorem{definition}[satz]{Definition}
\newtheorem{lemma}[satz]{Lemma}
\newtheorem{bemerkung}[satz]{Bemerkung}
\newtheorem{beispiel}[satz]{Beispiel}
\newtheorem*{bsp}{Beispiel}
\newtheorem*{bem}{Bemerkung}
\newtheorem*{defn}{Definition}
\newtheorem*{regeln}{Rechenregeln}

\makeatletter
\renewenvironment{proof}[1][\proofname]{\par\pushQED{\qed}%
  \normalfont\topsep\medskipamount\relax
  \trivlist\item[\hskip\labelsep\bfseries #1\@addpunct{.}]\ignorespaces}%
  {\popQED\endtrivlist\@endpefalse}
\makeatother
\renewcommand{\proofname}{Beweis}
\renewcommand{\qedsymbol}{$\square$}

\newcommand{\N}{\mathbb{N}}
\newcommand{\Z}{\mathbb{Z}}
\newcommand{\Q}{\mathbb{Q}}
\newcommand{\R}{\mathbb{R}}
\newcommand{\C}{\mathbb{C}}

\title{Analysis 1}
\author{Mitschrift}
\date{Wintersemester 2026/27}

\begin{document}

\maketitle
\tableofcontents
\newpage

% ======================================================================
% Kapitel 1
% ======================================================================
\section{Reelle und komplexe Zahlen}

\subsection{Mengen und Abbildungen}

Die Mengenlehre wird in dieser Vorlesung nicht axiomatisch konstruiert, sondern nur beschrieben. Für präzise Details (wie im Vorkurs) siehe die Literatur auf der Vorkurs-Webseite.

\begin{definition}[Naive Definition einer Menge, Cantor 1874]
Eine \emph{Menge} ist eine Zusammenfassung bestimmter, wohlunterschiedener Objekte unserer Anschauung oder unseres Denkens -- welche die \emph{Elemente} der Menge genannt werden -- zu einem Ganzen.
\end{definition}

\begin{beispiel}
\leavevmode
\begin{itemize}
  \item $M = \{2, 4, \pi, e\}$.
  \item $\N = \{0, 1, 2, 3, 4, \dots\}$ ist die Menge der \emph{natürlichen Zahlen}. \\
        (In dieser Vorlesung gehört die $0$ zu $\N$.)
  \item $\Z = \{\dots, -3, -2, -1, 0, 1, 2, 3, \dots\}$ ist die Menge der \emph{ganzen Zahlen}.
  \item $\Q = \left\{\tfrac{a}{b} \;\middle|\; a, b \in \Z,\ b \neq 0\right\}$ ist die Menge der \emph{rationalen Zahlen}, \\
        z.\,B.\ $1, -1, \tfrac12, -\tfrac12, \tfrac13, \tfrac23, \tfrac43 \in \Q$.
\end{itemize}
\end{beispiel}

\begin{bem}
Die Konstruktion dieser Zahlbereiche und der Rechenoperationen ($+,\ -,\ \cdot,\ :$) wird in dieser Vorlesung ohne Beweis vorausgesetzt und im 4.~Semester nachgeholt.
\end{bem}

\begin{bem}
Die obige naive Definition ist sehr schwammig. Viel schlimmer: Sie widerspricht sich selbst (Russellsche Antinomie, siehe Lineare Algebra~1 bzw.\ Vorkurs).

Nach Definition~1.1 wird eine Menge eindeutig durch die Elemente charakterisiert, die sie enthält. Z.\,B.\ wird die Menge $M = \{5, 7, 11\}$ eindeutig dadurch bestimmt, dass
\[
  5 \in M,\quad 7 \in M,\quad 11 \in M
\]
gilt und für alle anderen Objekte $x$ gilt: $x \notin M$.
\end{bem}

\begin{definition}
\leavevmode
\begin{enumerate}[label=(\alph*)]
  \item Die \emph{leere Menge} $\emptyset$ sei diejenige Menge, die kein Element enthält, d.\,h.\ für alle Objekte $x$ gilt $x \notin \emptyset$.
  \item Eine Menge $N$ heißt \emph{Teilmenge} einer Menge $M$, geschrieben $N \subseteq M$, falls jedes Element von $N$ auch Element von $M$ ist.
\end{enumerate}
\end{definition}

\begin{bsp}
$\{2, 4, 6\} \subseteq \{1, 2, 3, 4, 5, 6, 7\}$. Insbesondere ist $\emptyset \subseteq M$ für jede Menge $M$.
\end{bsp}

\paragraph{Notation.}
Teilmengen, die durch eine Bedingung beschrieben werden, notiert man in der Form
\[
  \{0, 1, 2, 3\} = \{\, n \in \N \mid n \leq 3 \,\},
\]
wobei \glqq$\mid$\grqq{} als \glqq mit\grqq{} oder \glqq für die gilt\grqq{} gelesen wird. Ein weiteres Beispiel:
\[
  \{4, 5, 6, 7\} = \{\, n \in \N \mid n > 3 \text{ und } n \leq 7 \,\}.
\]

\paragraph{Venn-Diagramme.}
In einem \emph{Venn-Diagramm} wird jede betrachtete Menge durch eine sie berandende Kurve dargestellt. Bei $n$ Mengen muss im Diagramm jede der $2^n$ Möglichkeiten für Enthaltensein bzw.\ Nicht-Enthaltensein eines Elements in den $n$ Mengen als Gebiet vorkommen.

\begin{figure}[H]
\centering
\begin{tikzpicture}[scale=0.85, every node/.style={font=\small}]
  % n = 2
  \draw[thick] (0,0) circle (1.2);
  \draw[thick] (1.2,0) circle (1.2);
  \node at (-1.15,1.2) {$N$};
  \node at (2.35,1.2) {$M$};
  \node at (-0.5,0) {\scriptsize $x \in N$};
  \node at (1.7,0) {\scriptsize $x \in M$};
  \draw[->] (0.6,1.9) node[above] {$x \in M \cap N$} -- (0.6,0.15);
  \fill (3.0,-1.1) circle (1.5pt) node[right] {$x \notin N \cup M$};
  \node at (0.6,-1.8) {$n = 2$};

  % n = 3
  \begin{scope}[xshift=7cm]
    \draw[thick] (-0.55,0.3) circle (1);
    \draw[thick] (0.55,0.3) circle (1);
    \draw[thick] (0,-0.6) circle (1);
    \node at (0,-1.8) {$n = 3$};
  \end{scope}
\end{tikzpicture}

\vspace{1em}
\begin{tikzpicture}[scale=0.75]
  % n = 4: vier Ellipsen, alle 16 Gebiete kommen vor
  \draw[thick, rotate around={-45:(-0.9,0)}] (-0.9,0) ellipse (2.4 and 1.15);
  \draw[thick, rotate around={-45:(0,0.8)}] (0,0.8) ellipse (2.4 and 1.15);
  \draw[thick, rotate around={45:(0,0.8)}] (0,0.8) ellipse (2.4 and 1.15);
  \draw[thick, rotate around={45:(0.9,0)}] (0.9,0) ellipse (2.4 and 1.15);
  \node at (0,-2.5) {$n = 4$};
\end{tikzpicture}
\caption{Venn-Diagramme für $n = 2, 3, 4$ Mengen}
\end{figure}

\begin{defn}
Für eine Menge $\mathcal{N}$ von Mengen ist die \emph{Vereinigungsmenge}
\[
  \bigcup_{M \in \mathcal{N}} M = \{\, x \mid \text{es gibt ein } M \in \mathcal{N} \text{ mit } x \in M \,\}
\]
und die \emph{Schnittmenge} (oder der \emph{Durchschnitt})
\[
  \bigcap_{M \in \mathcal{N}} M = \Bigl\{\, x \in \bigcup_{M \in \mathcal{N}} M \;\Bigm|\; \text{für alle } M \in \mathcal{N} \text{ gilt } x \in M \,\Bigr\}.
\]
\end{defn}

\begin{bsp}
Für $\mathcal{N} = \bigl\{\{1, 3, 5\}, \{3, 5, 7\}, \{3, 9, 11\}\bigr\}$ ist
\[
  \bigcup_{M \in \mathcal{N}} M = \{1, 3, 5, 7, 9, 11\}, \qquad \bigcap_{M \in \mathcal{N}} M = \{3\}.
\]
\end{bsp}

Für endlich viele Mengen $M_1, \dots, M_n$ schreibt man Vereinigung und Durchschnitt auch als
\[
  M_1 \cup M_2 \cup \dots \cup M_n \qquad\text{bzw.}\qquad M_1 \cap M_2 \cap \dots \cap M_n,
\]
z.\,B.\ $\{1, 3, 5\} \cap \{3, 5, 7\} \cap \{3, 9, 11\} = \{3\}$.

\begin{defn}
Für zwei Mengen $M$ und $N$ sei das \emph{Komplement} von $N$ in $M$ (die \emph{Differenzmenge})
\[
  M \setminus N = \{\, x \in M \mid x \notin N \,\}.
\]
\end{defn}

\begin{bsp}
$\{1, 3, 5\} \setminus \{3, 5, 7\} = \{1\}$.
\end{bsp}

\begin{figure}[H]
\centering
\begin{tikzpicture}[scale=0.85]
  \begin{scope}[even odd rule]
    \clip (1.2,0) circle (1.2) (0,0) circle (1.2);
    \clip (1.2,0) circle (1.2);
    \fill[pattern=north east lines] (1.2,0) circle (1.2);
  \end{scope}
  \draw[thick] (0,0) circle (1.2);
  \draw[thick] (1.2,0) circle (1.2);
  \node at (-1.15,1.2) {$N$};
  \node at (2.35,1.2) {$M$};
  \draw[->] (3.4,-1.2) node[right] {$M \setminus N$} -- (1.9,-0.3);
\end{tikzpicture}
\caption{Die Differenzmenge $M \setminus N$ (schraffiert)}
\end{figure}

\begin{regeln}
Die Mengenoperationen $\cup$, $\cap$ und $\setminus$ genügen etlichen Rechenregeln, z.\,B.\ gilt für Mengen $A, B, C$:
\begin{enumerate}[label=(\arabic*)]
  \item \emph{Kommutativgesetz:} \\
        $A \cup B = B \cup A$ \ und \ $A \cap B = B \cap A$,
  \item \emph{Assoziativgesetz:} \\
        $A \cup (B \cup C) = (A \cup B) \cup C$ \ und \ $A \cap (B \cap C) = (A \cap B) \cap C$,
  \item \emph{Distributivgesetz:} \\
        $A \cap (B \cup C) = (A \cap B) \cup (A \cap C)$ \ und \ $A \cup (B \cap C) = (A \cup B) \cap (A \cup C)$.
\end{enumerate}
\end{regeln}

\begin{proof}
Exemplarisch für $A \cup (B \cap C) = (A \cup B) \cap (A \cup C)$.

\glqq$\subseteq$\grqq: Sei $x \in A \cup (B \cap C)$. Ist $x \in A$, so ist $x \in A \cup B$ und $x \in A \cup C$, also $x \in (A \cup B) \cap (A \cup C)$. Ist $x \in B \cap C$, so ist $x \in B \subseteq A \cup B$ und $x \in C \subseteq A \cup C$, also ebenfalls $x \in (A \cup B) \cap (A \cup C)$.

\glqq$\supseteq$\grqq: Sei $x \in (A \cup B) \cap (A \cup C)$. Ist $x \in A$, so ist $x \in A \cup (B \cap C)$. Ist $x \notin A$, so folgt aus $x \in A \cup B$, dass $x \in B$, und aus $x \in A \cup C$, dass $x \in C$. Also $x \in B \cap C \subseteq A \cup (B \cap C)$.

Beide Inklusionen zusammen ergeben die Gleichheit.
\end{proof}

\paragraph{Tupel und kartesisches Produkt.}
Zu Mengen $M_1, \dots, M_n$ und Elementen $x_1 \in M_1, \dots, x_n \in M_n$ sei $(x_1, \dots, x_n)$ das \emph{geordnete $n$-Tupel} von $x_1, \dots, x_n$ bzw.\ $(x_1, x_2)$ das \emph{geordnete Paar}. Im Gegensatz zu den Elementen einer Menge $\{x_1, x_2, \dots\}$ ist dabei die Reihenfolge wichtig: Falls $x_1 \neq x_2$, dann ist
\[
  (x_1, x_2) \neq (x_2, x_1), \qquad\text{wohingegen}\qquad \{x_1, x_2\} = \{x_2, x_1\}.
\]

\begin{defn}
Die Menge aller $n$-Tupel ist das \emph{kartesische Produkt}
\[
  M_1 \times M_2 \times \dots \times M_n = \{\, (x_1, x_2, \dots, x_n) \mid x_1 \in M_1, \dots, x_n \in M_n \,\}.
\]
\end{defn}

\begin{bsp}
$\{1, 2\} \times \{3, 4, 5\} = \{(1, 3), (2, 3), (1, 4), (2, 4), (1, 5), (2, 5)\}$.
\end{bsp}

\paragraph{Die natürlichen Zahlen.}
Die Menge $\N$ der natürlichen Zahlen genügt dem folgenden Prinzip.

\begin{satz}[Prinzip der vollständigen Induktion]
Sei $A \subseteq \N$ eine Teilmenge mit
\begin{enumerate}[label=(\arabic*)]
  \item \emph{Induktionsanfang:} $0 \in A$,
  \item \emph{Induktionsschritt:} Für jedes $n \in A$ \textup{(}\emph{Induktionsannahme}\textup{)} ist auch $n + 1 \in A$.
\end{enumerate}
Dann folgt $A = \N$.
\end{satz}

\begin{bem}
Genauso kann man die Induktion statt bei $0$ auch bei jeder anderen Zahl $k \in \N$ starten (Induktionsanfang $k \in A$) und erhält dann die Aussage
\[
  \{\, m \in \N \mid m \geq k \,\} \subseteq A.
\]
\end{bem}

\begin{bem}
Ab hier beginnt der mathematisch exakte Teil der Vorlesung.
\end{bem}

Als Beispiel für einen Induktionsbeweis zeigen wir:

\begin{lemma}
Für endliche Mengen $M = \{x_1, \dots, x_m\}$ und $N = \{y_1, \dots, y_n\}$ mit $m$ bzw.\ $n$ Elementen hat $M \times N$ genau $m \cdot n$ Elemente.
\end{lemma}

\begin{proof}
Mit Induktion über $n$.

\emph{Induktionsanfang} ($n = 0$): Es ist $M \times \emptyset = \emptyset$, hat also $m \cdot 0 = 0$ Elemente.

\emph{Induktionsschritt} ($n \to n+1$): Angenommen, die Behauptung stimmt für ein bestimmtes festes $n$, d.\,h.\ $M \times \{y_1, \dots, y_n\}$ hat $m \cdot n$ Elemente. Dann besteht $M \times \{y_1, \dots, y_n, y_{n+1}\}$ aus diesen Paaren und außerdem den $m$ Paaren
\[
  (x_1, y_{n+1}), \dots, (x_m, y_{n+1}),
\]
hat also insgesamt $m \cdot n + m = m \cdot (n+1)$ Elemente. Somit stimmt die Behauptung auch für $n + 1$, nach dem Induktionsprinzip also für alle $n \in \N$.
\end{proof}

\end{document}
